\section{Galois preconditions and norm facts on the good cut}

% BEGIN C84F1
\paragraph{C84F1. Residual data and good multiplicities on the cut.}
Fix one datum $D$, norm choice $n$ and good cut $c$. At a retained
localization write $X$ for the old finite set of Hecke points,
$S\subset X$ for the nonempty retained set, $O$ for its coefficient
DVR, and $T\subset O^X$ for the original local Hecke algebra.
The cut algebra is its image in $O^S$. For any $f\in T$ its values
at all points have identical residue. Indeed the evaluation at
each point carries the constants onto $O$. The inverse image of
the maximal ideal of $O$ is therefore a maximal ideal of $T$.
Locality makes these ideals equal, and $f-f(x)$ is in the one at
$x$. It follows coefficient by coefficient that the reduced Hecke
polynomials at $x,y$ agree at every test prime.

A residual rank-$r$ representation is tested by its continuity
condition and its Frobenius characteristic polynomials. A test at
every retained point implies the test at an arbitrary old point:
choose one retained point and use the equality of reduced Hecke
polynomials just proved. Restriction gives the converse. The same
argument with the same residue polynomial witnesses transports
$\mathrm{AtM}$, the condition that a given automorphic label occur
in the residual localization. Thus the old and new good parts in
the complex source are the same subspace.

For $x\in S$ a representative has new point $x$ exactly when its
old point is the underlying $x$. Hence the possible parameter
lists and their complex type spaces are unchanged, as are the
sets $Z$ and $Z_{\rm good}$ after intersection with $S$.
Goodness gives the reverse inclusion $Z_{\rm good}\subset S$.
Consequently the sum of the indicated complex spaces over good
points is the same. T83C2 identifies the full rational eigenspace
at a retained point isomorphically with the original one. Its
dimension, the local multiplicity, is the same; excluded points
make no contribution to the sum weighted by goodness. These
facts transfer each image, kernel and dimension identity of
$\mathrm{ConclusionZ}$, the good-part clause in row C, and its
good-point nonvanishing test. Local uniqueness transfers because
both old $\mathrm{AtM}$ predicates hold whenever their cut versions
hold, and equality of two underlying retained indices implies
their equality as indices of the cut. The residual irreducibility
and tame-parameter tests transfer by the residual and parameter
identifications above.
% END C84F1

% BEGIN C84F2
\paragraph{C84F2. The integral coefficient and row data of the cut.}
Assume the original $\mathrm{GaloisPreconditions}(D,n)$.
A live cut index is live on $D$: its witness is the underlying
point, with the same $Z$ label and local prime. Take at that index
the original smaller coefficient DVR $A$, its embedded Hecke
algebra $T_A$, the integral models of the two representations and
the characters. On $S$ use the image of $T_A$ by restriction and
the restrictions of the same matrices. The $O$-span of their
embedded Hecke algebra is the new $O$-Hecke algebra: applying the
linear restriction map to $O\cdot T_A=T_O$ commutes with sums and
scalar multiplication. Its image is local, and the restriction
of a separator for any two retained points still separates them.
Equations between characteristic polynomials, the chosen changes
of basis and the projected idempotents hold after evaluation at
every point of $S$. For a matrix in the algebra generated by the
representation use its finite expression $\sum f_j r(g_j)$,
restricting each coefficient. In the congruence module the
additional coefficients must vanish at the selected point; they
retain this property after restriction. This transports the
chosen congruence identities integrally, including their original
uniformizer power. These witnesses give the full cut
$\mathrm{Large}$ condition.

The global fields, infinity types, and analytic and geometric
coefficient embeddings of the datum are the same. The
automorphic Galois admissibility test at a retained index uses
the same semisimple realization, with all its local equalities,
its inertia and monodromy comparisons and its de Rham tests.
The signature and temperedness tests therefore restrict at once,
as do the weight and block-character dictionaries, unramified
twisting and depth-zero tests. Combined with C84F1 this gives
the cut admissible datum and its rows A, B, C. The vacuous
auxiliary slot in row A has predicate $\mathrm{True}$. For the
rank test on a size datum on the cut use its identical scale
and numerical fields on the original member. This gives the
original rank bound for precisely that cut input.

Here is a way to transfer the first provision without rerunning
the automorphic construction. In its family clauses a cut
$\mathrm{GoodR}$ can be extended by the given original good
integral realization at every discarded point: its defining
Hecke, continuity and semisimplicity tests are pointwise. All
retained characteristic-polynomial evaluations and parameters
are the old ones. For the clause that characteristic coefficients
belong to the Hecke algebra, take the restriction of the old
coefficient in $T$. For the residual comparisons use C84F1.
For tests conditional on liveness, a live cut index is live
upstairs. For pointwise tests at an unretained automorphic
representative the hypothesis that it has a specified cut point
forces its underlying point to be retained. The family of
rational realizations is obtained by restriction as well.
For the parameter clause $42$ the hypothesis is a nonzero
isotypic row. Its $\theta$-images, its type subspaces and the
definition of a parameter of that row all use the original
global row and automorphic parameter spaces at the underlying
index. They do not use the point-valued option of a
representative. The old positive scalar, theta representative
and parameter list thus serve at the cut, as does cuspidality
on a good row. This completes the first provision.

The cut norm choice is $n_c=\mathrm{nchCut}(c,n)$. Its twisted
place, side, auxiliary set and character ties restrict to the
underlying index, with the old common character witness. The
triple, equal-twist, Fourier-type, inert base-change and
$\lambda$ tests retain their original witnesses. Together with
C84F1 this gives $\mathrm{GaloisPreconditions}(D_c,n_c)$,
where $D_c=\operatorname{cut}(D,c)$. Also $D_c$ has
$\mathrm{CutRec}$: any representative with a cut point has an
old point in $S$, so the nice-local-parameter clause of goodness
applies to that very representative at both generator primes.
% END C84F2

% BEGIN C84F3
\paragraph{C84F3. All labelled de Rham jumps and the admissible jump datum.}
In this and the next three paragraphs work on any $D$ with
$\mathrm{GaloisPreconditions}(D,n)$ and $\mathrm{CutRec}(D)$.
Fix $i,r,x$ with $r$ good and a place $w\mid\ell_i$. Choose an
automorphic representative $P$ over $x$ and a base-change list
$\mathcal P$. At a generator prime, $w$ splits over the real
subfield and its local parameter for $P$ equals that of
$\mathcal P$. The cut hypothesis gives the prescribed tame,
zero-monodromy and inertia-root multiplicity condition $\mathrm{Nc}$ on
this parameter. At any other coefficient prime the parameter
is unramified: at a split place this is the unramified level
condition off the two generator primes and $S_i,Q_i$, and at
an inert place it is in the base-change definition. An
unramified parameter satisfies $\mathrm{Nc}$.

Form the semisimple Galois realization associated to
$\mathcal P$ using automorphic Galois admissibility on its
constituent parts and their prescribed Tate shifts. Its split
Frobenius polynomials agree with those of $r_x$. Density and
semisimplicity give a conjugacy over $\overline{\mathbb Q}_{\ell_i}$.
The admissible comparison for $\mathrm{Nc}$ gives de Rhamness and
the labelled Hodge weights of this realization, and the
conjugacy clause of the comparison transfers both to $r_x$ at
each embedding inducing $w$. Thus its jump multiset is the
prescribed level-zero multiset. Away from the generator primes
the crystalline assertion and the same weight equality are
already in the first provision. Moreover the weights in this
multiset are distinct by regularity of the original infinity
types. If a nonzero irreducible constituent appeared twice
after splitting, de Rhamness of summands and additivity for
direct sums would repeat one of its jumps at any inducing
embedding. This proves distinctness. A common finite splitting
extension for the finite set of points suffices: choose bases of
the absolutely irreducible summands of their semisimple
representations over the algebraic closure and adjoin all
entries of the change-of-basis matrices. This proves the two
native $g2$ assertions.

For $\mathrm{admJ}$ take an arbitrary size datum $L$ with
$\mathrm{AdmPar}(L)$. For a de Rham representation the member's
Weil--Deligne comparison identifies crystallinity with unramified
Weil--Deligne parameter. Conversely a crystalline representation
is de Rham by the crystalline/de Rham period inclusion. These
comparisons may be tested over a finite coefficient extension:
the inverse images of $\mathrm{GL}_a$ of the finite extensions
cover the compact Galois group by countably many closed subgroups.
One is open by Baire, and adjoining the matrix entries of its
finitely many coset representatives suffices. Entrywise
restriction to that finite field preserves continuity.
For the weight dictionary distinguish the chosen archimedean
embedding over the real place and its conjugate. The CM type
contains exactly one of the two. Use the member's weight
parameter at that embedding and its block character there; the
common finite-order factors do not change the labelled jumps of
the character. The arithmetic character of the Galois realization
and the chosen automorphic character agree at the split
Frobenius tests and hence on the Galois group by density.
This gives the two signs and the common shift of the native
dictionary (the global orientation parameter is false). More
explicitly, with $\sigma=\operatorname{sgnOf}(\mathrm{true},b)$
for the embedding choice $b$, the level-zero multiset is
$\{\sigma u+c_\tau:u\in \mathrm{PP}(D,i,\rho)\}$ with
multiplicities, and the first character has its single jump at
$-\sigma e_\rho(\mathrm{liftIdx}(i))-(t-1)/2+c_\tau$.
In the local reading, compose the given embedding of the
completion with $E\to E_w$. It induces $w$. Under a rank-one
twist with labelled weight $c$ the degrees shift by that
weight, and duality reverses the grading. With the prescribed
Tate twist the reflected label is $c-1-a$ at degree $a$.
Additivity on a given direct-sum decomposition and the test
$\mathrm{gr}^a\ne0$ if and only if its dimension is nonzero
give the ordinary and dual summand tests over every admitted
coefficient extension. These are the $\mathrm{AdmRead}$ clauses;
the two $g2$ conclusions above and the given $\mathrm{AdmPar}$
supply the remaining entries of $\mathrm{AdmJAlt}(L)$.
% END C84F3

% BEGIN C84F4
\paragraph{C84F4. The tame polynomial and vanishing of the local cohomology in G5.}
Use local finiteness, local Tate duality and the local Euler
characteristic theorem. At an index satisfying $\mathrm{HG}$ choose
a good point and its good parameter. Row C supplies a residual
rank-$r$ realization there. Its Frobenius tests at that point
are those at every point by C84F1's local Hecke calculation.
We may thus choose a global residual part $\bar\rho$. For any
such part, every good family $r$, and every $w\in S_i\cup Q_i$,
the base-change comparison at $w$ identifies inertia polynomials
and gives zero monodromy. Indeed $w\nmid\ell_i$ and splits.
Use the same semisimple conjugacy as in C84F3. Its
Weil--Deligne relation at $w$ gives the full dimension of the
local parameter. At a generator place the needed $\mathrm{Nc}$
comes from the cut; at $Q_i$ and the specified auxiliary primes
the given monodromy/inertia type and its full dimension give
$\mathrm{Nc}$. In the split case of the G5 statement that type
has $r$ trivial characters and $t$ copies of $\chi_1$. Vanishing
monodromy gives a finite inertial image; the simultaneous
character splitting is over the fraction field because both
characters are defined there. Consequently
$(r_x(\gamma')-1)(r_x(\gamma)-\chi_1(\gamma))=0$ for every
two inertia elements there.

Set $z=\overline{\chi_1(\gamma_0)}$ for the prescribed tame
generator at $w\in S_i$. At either generator prime $z$ is excluded
from the roots of $\bar\rho(\gamma_0)$ by residual admissibility:
$w\nmid\ell_i$ supplies its unequal-prime guard. At the remaining
auxiliary prime, in the split calculation above the full residual
polynomial is $(X-1)^r(X-z)^t$. The global residual
factorization is $\operatorname{char}_{\bar\rho}(X)(X-z)^t$
because the cyclotomic character is trivial on this inertia.
Cancel the monic factor. Here the auxiliary inertia character
has image of order $15$ and the designated generator has
primitive value. Indeed wild inertia at this prime, whose
residue characteristic is $1$ modulo $15$, maps trivially to
the cyclic image of order $15$. A primitive fifteenth root has
unit difference from one since $\Phi_{15}(1)=1$. Thus in either
case $z$ is a nonroot. For every integer $e$ the residual twist
$M=\bar\rho\otimes\bar\chi_1^{-1}\bar\epsilon^{e}$ has
$\gamma_0-1$ invertible. This kills its invariants. The same is
true on its Tate dual: on inertia the cyclotomic character is
trivial and inverse transpose preserves the absence of eigenvalue
one. Local Tate duality kills $H^2$; the finite local Euler
characteristic away from $\ell_i$ gives
$|H^1|=|H^0||H^2|=1$. These are all the required twists.

At a generator place in $S_i$ form the polynomial of
$r_x(\gamma_0)$ simultaneously in $T[X]$ by the coefficient
clause of the first provision. Its reduction has root $z$ of
multiplicity exactly $t$ by the nonroot calculation. Its
characteristic-zero multiplicity at $\chi_1(\gamma_0)$ is at
most $t$, since reduction can only increase this multiplicity.
Apply the generator-prime clause of the input
$\mathrm{MonodromyInertiaType}$ at $w\in S_i$ over $g(k)$:
its two tests are precisely that $\gamma_0$ generates tame
inertia and this characteristic-zero multiplicity is at most
$t$. It supplies a decomposition of the local parameter into
parts of dimensions $r,t$. The second part has zero monodromy
and characteristic polynomial $(X-\chi_1(\gamma))^t$ for
every inertia element $\gamma$. The Weil--Deligne comparison
for a direct sum gives complementary inertia-stable spaces.
On the second one finite inertia acts semisimply in
characteristic zero, hence every $\gamma$ acts as the scalar
$\chi_1(\gamma)$. Divide the universal characteristic polynomial by the
monic $(X-\chi_1(\gamma_0))^t$ in $T[X]$. Its remainder is zero
at all $x$, hence zero in $T[X]$. The quotient $P$ is the
characteristic polynomial of the complementary block at each
$x$. Cayley--Hamilton there, and the character action on the
other block, give
$P(r_x(\gamma_0))(r_x(\gamma)-\chi_1(\gamma))=0$ integrally.
Its value at $\chi_1(\gamma_0)$ is a unit in the local algebra
by the residual multiplicity. This supplies one polynomial for
all points and all inertia elements.

Finally apply the idempotent-choice functional of G5 to the
corresponding matrix relation on its congruence generator
$y_\gamma$, a right multiple of $r(\gamma)-\chi_1(\gamma)$
with the prescribed scalar normalization. The left operator
$P(r(\gamma_0))$ belongs to the indicated Hecke span. Its
character acts on this functional by the displayed unit, so
its vanishing on the product forces vanishing on $y_\gamma$.
In the auxiliary case use $r(\gamma')-1$ in place of the
polynomial. Choose $\gamma'$ giving a primitive fifteenth root;
its character minus one is a unit and the two-element split
identity above gives the zero product. This proves all clauses
of $\mathrm{ConclusionG5}$, including the integral functional
test.
% END C84F4

% BEGIN C84F5
\paragraph{C84F5. A character eigenline detects the full rank-one string in G6.}
Take a good index, a good integral family $r$, and a point $x$ at
which $r_x$ has a nonzero $\chi_1$ eigenvector over the fraction
field. Choose a representative and any one of its base-change
lists $(a_k,\pi_k,d_k)_k$. For the $k$th part take the admissible
centered Galois realization $R_k$ of rank $a_k$. For the required
algebraic twist, the top weight of this string at infinity is the
weight of $\pi_k$ translated by $(d_k-1)/2$. Its multiset is
contained in the regular algebraic weight of the entire list, so
it has no repeats and lies in the required translated integer
lattice. This is the parity and regularity condition for taking
the centered twist of the attached cuspidal realization in the
given automorphic Galois admissibility hypothesis. The polynomial
of $r_x(g)$ is that of the direct sum
\[
               \bigoplus_k\bigoplus_{0\le j<d_k}
                                     R_k\epsilon^{-j}.
\]
Indeed both polynomials at every split unramified Frobenius away
from the finitely many excluded places are the factors of the
same base-change parameter, with $\epsilon(\mathrm{Frob}_w)=q_w$.
All coefficients are continuous, and these Frobenius conjugacy
classes are dense. Both sides are semisimple over the algebraic
closure. Brauer--Nesbitt and projection of the nonzero eigenline
therefore give a $\chi_1$-eigenvector in at least one shifted
summand, say $(k,j)$.

At a split place at which $\pi_k$ and the distinguished character
are unramified, the vector gives eigenvalue
$q_w^{j}\iota(\chi_1(\mathrm{Frob}_w))$ in $R_k$. Compare its
absolute value with the normalization of the centered parameter.
For the corresponding unshifted cuspidal Satake root $\alpha$,
this gives $2\bigl|\log_{q_w}\lvert\alpha\rvert\bigr|=|2j+t-d_k|$. The native open cuspidal bound confines its absolute
value strictly between $q_w^{-1/2}$ and $q_w^{1/2}$.
The bound applies also when the integral parity condition for the
cuspidal part requires the fixed half-weight character: use the
admissible twist and remove its known norm in the Satake and
infinity parameters. Since $2j+t-d_k$ is an integer and $q_w>1$,
we obtain $d_k=t+2j$.

For each $k$ let $A_k$ be the multiset of dimensions in the
residual semisimplification of a stable lattice of $R_k$, after
extending coefficients to a common algebraic closure of the
residue field. Such lattices exist: the matrices descend to a
finite local coefficient field as in C84F3, and the span of the
translates of a lattice lies in a bounded multiple of it by
compactness. Each $A_k$ has positive entries of sum $a_k$.
Reduction of the polynomial identity gives
\[
            \sum_k d_k A_k=\{r\}+t\{1\}.
\]
In fact the comparison on the right is the residual polynomial
of C84F4's absolutely irreducible rank-$r$ part and the $t$
characters. Semisimplify after reduction on both sides; twisting
preserves irreducibility and dimension. The polynomial at every
group element determines the semisimple constituents: the
residual lists factor through a common finite quotient, and on
its prime-to-$\ell$ elements the matching eigenvalue multisets
give the same Brauer character. Taking dimensions with
multiplicities gives the stated equality.

Here $r>1$ and $t\ge2$. Our selected $k$ has $d_k\ge t$.
An entry larger than one in $A_k$ would occur at least twice
on the left and at most once on the right. Every entry is
therefore one, and there are at most $t$ available copies,
forcing $A_k=\{1\}$, $a_k=1$, $d_k=t$ and $j=0$.
The eigenline is now the whole $R_k=\chi_1$. The prescribed
one-dimensional automorphic character $\mu_i\eta_T^{-1}$
has the same normalized realization: its Frobenius values
agree with $\chi_1$ by the given character tie, hence agree
on the whole group by density. Undoing the common local twist
identifies the corresponding local automorphic characters almost
everywhere; admissible multiplicity one identifies them globally.
This identifies the required part in the base-change list, so
$x\in Z_N$.

Finally every parameter $\Sigma$ at $x$ gives another such list
by appending the prescribed rank-one part and twisting. Apply
the residual multiset calculation to this list. Its appended
part uses all $t$ copies of $1$, leaving the singleton $\{r\}$.
Every remaining part has nonempty $A_k$ and $d_k\ge1$. Exactly
one remains, with $d_k=1$ and $a_k=r$. Undoing the twist gives
$\Sigma=[(r,\pi,1)]$. This proves both native conclusions for
every tested eigenline, including the universal parameter
quantifier.
% END C84F5

% BEGIN C84F6
\paragraph{C84F6. Coefficient duties and the actual cut norm-choice facts.}
The four coefficient duties use the same residual representation,
ambient point family and smaller coefficient model as above.
The first is $\mathrm{OLarge}$ from $\mathrm{Large}$. For distinct
constituents, use C84F3 on a finite splitting field. The
$\mathrm{OSplit}$ clause of $\mathrm{Large}$ splits the original
semisimple good representation into absolutely irreducible pieces
over the original fraction field. If two disjoint pieces there
were isomorphic, they would remain so after scalar extension,
contradicting C84F3. This supplies the required condition over
the original coefficients.

For the design test the given residual matrix span on the
rank-$r$ summand is the full matrix algebra; after any residue
extension a stable subspace for all these matrices is either
zero or the whole space, by the elementary matrices. This is
absolute irreducibility. At a type-prime test the given rank
decomposition instead identifies the residual part by dividing
the Hecke characteristic polynomial by the prescribed character
factors at split Frobenius places. C84F1's residual comparison
makes this a residual part at every point. The original
residual irreducibility test now supplies the same conclusion.
In either case the chosen point and the given decomposition
witness liveness, so $\mathrm{Large}$ supplies its smaller-ring
model including the projected block $e_B$, the complementary
rank-one matrix $e_1$, the regular element $X_0$ and the
specified $g_0$ in inertia.

Here are the integral issues in using that model for the
representation in the test. Extend the smaller ring into $O$;
the chosen block polynomial, monic of degree $t$, has $t$
distinct roots and already splits over the specified fraction
field. Its roots lie in $O$ by integrality. Over each point take
the $t!$ orderings of those roots. Over $T$ the successive monic
root adjunctions have bases of lengths $t,t-1,\ldots,1$.
Their product basis is independent in the ring of functions on
these ordered tuples: over each fraction-field point the roots
are distinct, giving the product of fields by successive
interpolation, and the original $T$ embeds in their product.
Thus the $T$-algebra generated by the root functions has its
required integral free basis, by the same monic-division argument
as D79G5, with distinctness needed only over the fraction field. At the distinguished lift order
them as $\chi_c(g_0)$, $1\le c\le t$. Multiplying the other
linear factors of $X_0$ in its $e_B$ block selects the first
line multiplied by the product of its root differences.
At every ordering this product of factors has rank one and its
square is its scalar root-product times itself, by evaluating on
each distinct eigenspace in the block. The characteristic and
block determinant coefficients belong to the Hecke span by the
word expansion of D79G2. These identities hold integrally by
injectivity into the fraction field. The residual
absolute irreducibility gives all integral rank-block matrices
as evaluations of the chosen algebra: reduction spans the
matrix algebra by the density theorem, and Nakayama lifts the
span. The resulting ordered-root record has the original
$g_0$ and the prescribed product in its length clause.

For a different tested good family, density gives identical
rational characteristic polynomials pointwise, hence conjugate
semisimple representations over the fraction field. Cancelling
the prescribed character summands gives the conjugacy of the
rank blocks at the tested point. Transport every matrix in the
generated algebra by evaluating its original Hecke-linear
group-algebra expression in the new integral matrices. This
defines both it and the inverse transport integrally because
conjugation over the fraction field detects equality of two
expressions. The rank-block action on its integral lattice
still gives every matrix: the conjugacy $Q$ on the rank block
carries every old elementary matrix into an integral matrix
(the conjugated expression acts integrally by a sum of the
new rank representation matrices). Thus $Q M_r(O)Q^{-1}$
lies in $M_r(O)$. Conjugation on the full matrix space has
determinant one; both orders are lattices of the same covolume,
so containment is equality. Character evaluations, ranks,
determinants and the ordered-root functions stay the same.
The element $g_0$ thereby retains its required wild cyclotomic
condition at a type prime. These are respectively the
ordinary and type-prime coefficient duties.

Apply C84F3--C84F6 to $(D_c,n_c)$ of C84F2. In the norm-choice
record, the second provision consists precisely of C84F4,
C84F5 and the good-point identities of C84F1. Insert the $g2$
and parameter conclusions of C84F3 and the coefficient duties
just proved. Every remaining entry is the corresponding cut
precondition of C84F2. This gives
$\mathrm{NormChoiceFacts}(D_c,n_c)$ for every good $c$ on a
member with the stated original preconditions. The arithmetic
principles used here are the given admissibility and row
hypotheses, Chebotarev, and local finiteness, Tate duality and
the local Euler characteristic formula. In particular D79G7
used G5 and G6 in its proof on a supplied hub. The two
computations above supply them in the reverse direction
required here from the preconditions and $\mathrm{CutRec}$.
% END C84F6
